\documentclass[11pt]{amsart}
\usepackage[margin=1.15in]{geometry}
\usepackage[T1]{fontenc}
\usepackage{lmodern,microtype}
\usepackage{amsmath,amssymb,amsthm,mathtools,booktabs}
\usepackage[colorlinks=true,linkcolor=blue,citecolor=blue,urlcolor=blue]{hyperref}

\newtheorem{theorem}{Theorem}[section]
\newtheorem{lemma}[theorem]{Lemma}
\newtheorem{proposition}[theorem]{Proposition}
\newtheorem{corollary}[theorem]{Corollary}
\newtheorem{conjecture}[theorem]{Conjecture}
\theoremstyle{definition}
\newtheorem{definition}[theorem]{Definition}
\theoremstyle{remark}
\newtheorem{remark}[theorem]{Remark}
\numberwithin{equation}{section}

\newcommand{\KG}{\operatorname{KG}}
\newcommand{\up}{\operatorname{up}}

\hypersetup{pdftitle={A characterization of 3-tight graphs},pdfauthor={Brandon Li}}

\title{A characterization of $3$-tight graphs}
\author{Brandon Li}
\address{Great Neck, New York, USA}
\email{brandon.li.gn@gmail.com}
\subjclass[2020]{05C15, 05D05}
\keywords{stable Kneser graph, chromatic number, intersecting family, Schrijver graph}

\begin{document}

\begin{abstract}
The $G$-stable Kneser graph $\KG(G,k)$ has the independent $k$-sets of a graph
$G$ as vertices, two of them adjacent when they are disjoint. If $G$ has $n$
vertices, then $\chi(\KG(G,k))\le n-2k+2$, and $G$ is called $k$-tight when
equality holds. Arana, Haxell, Kaiser and Stehl\'ik proved that in a $k$-tight
graph every component is a path or a cycle, and every cycle other than $G$
itself has length at least $4k-3$. They conjectured the converse. We prove the
conjecture for $k=3$: a graph on at least seven vertices is $3$-tight if and
only if it is a cycle or each of its components is a path or a cycle of length
at least nine. The proof combines a minimal counterexample, the Polcyn--Ruci\'nski
classification of maximal intersecting triple systems, and a triangle count,
which leave ten small graphs that are settled by exact computer checks. For
every $k\ge3$ we also show that a smallest counterexample has fewer than
$k^3-k$ vertices.
\end{abstract}

\maketitle

\section{Introduction}

Let $G$ be a graph on $n$ vertices and let $k\ge2$. The \emph{$G$-stable
Kneser graph} $\KG(G,k)$ has as its vertices the independent sets of size $k$
in $G$, and two of them are adjacent when they are disjoint. When $G$ has no
edges this is the Kneser graph, whose chromatic number is $n-2k+2$ by
Lov\'asz's theorem. When $G$ is the cycle $C_n$ it is the Schrijver graph, and
Schrijver~\cite{Schrijver} proved that its chromatic number is still $n-2k+2$.
Since $\KG(G,k)$ is an induced subgraph of the Kneser graph, $n-2k+2$ is an
upper bound for every $G$.

Arana, Haxell, Kaiser and Stehl\'ik~\cite{AHKS} asked which graphs attain this
bound. Following them, a graph $G$ on $n\ge2k+1$ vertices is \emph{$k$-tight}
if $\chi(\KG(G,k))=n-2k+2$. They proved the following necessary condition.

\begin{theorem}[{\cite[Theorem~1.2]{AHKS}}]\label{thm:necessary}
Let $k\ge2$. If a graph $G$ on at least $2k+1$ vertices has a vertex of degree
at least~$3$, or properly contains a cycle of length at most $4k-4$, then $G$
is not $k$-tight.
\end{theorem}

They conjectured that this condition is also sufficient.

\begin{conjecture}[{\cite[Conjecture~1.3]{AHKS}}]\label{conj:AHKS}
For any $k\ge2$, a graph $G$ on at least $2k+1$ vertices is $k$-tight if and
only if either $G$ is a cycle, or every component of $G$ is a path or a cycle
of length at least $4k-3$.
\end{conjecture}

For $k=2$ the conjecture holds: it is equivalent to a theorem of Daneshpajouh,
Meunier and Mizrahi~\cite{DMM}, as explained in~\cite[Theorem~1.1]{AHKS}.
For $k\ge3$, Arana, Haxell, Kaiser and Stehl\'ik point out that even the
simplest case not covered by Schrijver's theorem, a cycle of length at least
$4k-3$ together with one isolated vertex, is open, and they prove it for all
sufficiently long cycles~\cite[Theorem~1.4]{AHKS}. For $k=3$ this case was
settled for cycle lengths $9$, $10$ and at least $15$ in~\cite{AFTD}.

Our main result is the case $k=3$ of Conjecture~\ref{conj:AHKS}.

\begin{theorem}\label{thm:main}
A graph $G$ on $n\ge7$ vertices is $3$-tight if and only if $G$ is a cycle, or
every component of $G$ is a path or a cycle of length at least~$9$.
\end{theorem}

The forward implication is Theorem~\ref{thm:necessary}. For the converse we
study a smallest counterexample. Section~\ref{sec:prelim} shows that in a
smallest counterexample every color class of an optimal coloring is an
intersecting family with no common vertex, for every~$k$. For $k=3$, such
maximal families were classified by Polcyn and Ruci\'nski~\cite{PR}. In
Section~\ref{sec:large} we combine their classification with a count of
triangles in an auxiliary graph to show that a smallest counterexample has at
most $13$ vertices. This leaves ten graphs, each a cycle of length at least
nine together with a path. Section~\ref{sec:isolate} handles the four graphs
in which the path is a single vertex, and Section~\ref{sec:small} handles the
other six. Both use exact computer checks.

Theorem~\ref{thm:main} contains the case of a cycle with any number of
isolated vertices.

\begin{corollary}\label{cor:isolates}
Let $c\ge3$, $r\ge1$ and $c+r\ge7$. The graph $C_c\sqcup rK_1$ is $3$-tight if
and only if $c\ge9$.
\end{corollary}

\begin{proof}
If $c\ge9$, the graph satisfies the condition of Theorem~\ref{thm:main}. If
$c\le8$, it properly contains a cycle of length at most $8=4k-4$, so it is not
$3$-tight by Theorem~\ref{thm:necessary}.
\end{proof}

The method also gives a bound on the size of a counterexample for every rank
(Section~\ref{sec:allk}).

\begin{theorem}\label{thm:allk}
Let $k\ge3$. If Conjecture~\ref{conj:AHKS} fails for $k$, then it fails for
some graph on at most $k^3-k-1$ vertices.
\end{theorem}

For comparison, the one-isolate case of~\cite{AFTD} holds for cycles of length
at least $k((k-1)^2+1)$, which is also cubic in~$k$. Theorem~\ref{thm:allk}
applies to all graphs in the conjecture. The conjecture for $k\ge4$ remains
open.

\section{Preliminaries and minimal counterexamples}\label{sec:prelim}

All graphs are finite and simple. For a graph $G$ we write $V(G)$ and $E(G)$
for its vertex and edge sets. A set $S\subseteq V(G)$ is \emph{independent} if
it contains no edge of $G$. We write $\mathcal I_k(G)$ for the family of
independent $k$-sets of $G$, which is the vertex set of $\KG(G,k)$. A family of
sets is \emph{intersecting} if every two of its members meet. It is a
\emph{star} if all its members contain a common element, and a \emph{nonstar}
otherwise. A color class of a proper coloring of $\KG(G,k)$ is exactly an
intersecting subfamily of $\mathcal I_k(G)$.

\begin{definition}
Let $k\ge2$. A graph $G$ is \emph{admissible} (at rank $k$) if $G$ is a cycle,
or every component of $G$ is a path or a cycle of length at least $4k-3$.
Single vertices count as paths. A \emph{path forest} is a graph whose
components are all paths.
\end{definition}

\begin{lemma}\label{lem:upper}
If $G$ has $n\ge2k$ vertices, then $\chi(\KG(G,k))\le n-2k+2$.
\end{lemma}

\begin{proof}
Order the vertices as $v_1,\dots,v_n$. Give a set $S\in\mathcal I_k(G)$ the
color $i$ if its smallest index is $i\le n-2k+1$, and the color $n-2k+2$
otherwise. Sets of color $i\le n-2k+1$ share $v_i$. Sets of color $n-2k+2$ lie
in the last $2k-1$ vertices, so any two of them meet.
\end{proof}

\begin{lemma}\label{lem:spanning}
If $H$ is a spanning subgraph of $G$, then $\KG(G,k)$ is an induced subgraph
of $\KG(H,k)$, and so $\chi(\KG(G,k))\le\chi(\KG(H,k))$. In particular, every
path forest on $n\ge2k+1$ vertices is $k$-tight.
\end{lemma}

\begin{proof}
Every independent set of $G$ is independent in $H$, and adjacency in both
graphs means disjointness. A path forest on $n$ vertices is a spanning
subgraph of $C_n$, and $C_n$ is $k$-tight by Schrijver's
theorem~\cite{Schrijver}. Lemma~\ref{lem:upper} gives the upper bound.
\end{proof}

A \emph{counterexample} (at rank $k$) is an admissible graph on at least $2k+1$
vertices that is not $k$-tight. It is \emph{minimal} if no counterexample has
fewer vertices. Note that admissibility depends only on the components.

\begin{lemma}\label{lem:minimal}
Let $k\ge2$ and let $G$ be a minimal counterexample on $n$ vertices. Then the
following hold.
\begin{enumerate}
\item[(a)] $G$ has a cycle component of length at least $4k-3$ and at least
one other component. In particular $n\ge4k-2$ and $|E(G)|\le n$.
\item[(b)] $\KG(G,k)$ has a proper coloring with $n-2k+1$ colors, and in every
such coloring each color class is a nonstar.
\item[(c)] Let $G'$ be obtained from $G$ by adding edges that join all path
components of $G$ into a single path. Then $G'$ is also a minimal
counterexample, so \textup{(a)} and \textup{(b)} hold for $G'$.
\end{enumerate}
\end{lemma}

\begin{proof}
(a) By Schrijver's theorem and Lemma~\ref{lem:spanning}, $G$ is neither a
cycle nor a path forest. Since $G$ is admissible, it has a cycle component,
of length at least $4k-3$, and another component. Every vertex has degree at
most two, so $|E(G)|\le n$.

(b) Since $G$ is not $k$-tight, Lemma~\ref{lem:upper} gives a coloring with
$n-2k+1$ colors. Suppose some color class is a star with common vertex $v$.
Deleting $v$ from a path leaves at most two paths, and deleting it from a cycle
component of length at least $4k-3$ leaves a path. The other components are
unchanged, so $G-v$ is admissible. By~(a), $G-v$ has $n-1\ge2k+1$ vertices,
so by minimality it is $k$-tight: $\chi(\KG(G-v,k))=n-2k+1$. But
$\KG(G-v,k)$ is the induced subgraph of $\KG(G,k)$ on the sets avoiding $v$,
and none of these has the color of the star. So $\KG(G-v,k)$ is colored with
at most $n-2k$ colors, a contradiction.

(c) The components of $G'$ are the cycles of $G$ and one path, so $G'$ is
admissible. By Lemma~\ref{lem:spanning},
$\chi(\KG(G',k))\le\chi(\KG(G,k))<n-2k+2$. So $G'$ is a counterexample with
$n$ vertices, hence minimal.
\end{proof}

For $k=3$ and $n\le13$ this pins down the shape of a minimal counterexample:
two cycles of length at least $9$ need $18$ vertices, so by
Lemma~\ref{lem:minimal}(a) and~(c) we may take it to be $C_c\sqcup P_r$ with
$c\ge9$, $r\ge1$ and $c+r=n\le13$. Here $P_r$ is a path with $r$ vertices.

\section{A cycle and an isolated vertex}\label{sec:isolate}

In this section $k=3$. Let $c\ge9$ and let $H=C_c\sqcup K_1$, where the
isolated vertex is called $\infty$. Then $H$ has $c+1$ vertices, so $H$ is
$3$-tight if and only if $\KG(H,3)$ has no proper coloring with $c-4$ colors.
We say that a triple $U$ \emph{covers} a triple $W$ if $|U\cap W|\ge2$, that
is, if $W$ contains one of the three pairs in~$U$.

\begin{lemma}\label{lem:anchors}
Let $c\ge9$. If $\KG(C_c\sqcup K_1,3)$ has a proper coloring with $c-4$
colors, then there are independent triples $U_1,\dots,U_{c-4}$ of $C_c$ with
the following properties.
\begin{enumerate}
\item[(i)] No two of them share a pair.
\item[(ii)] Every independent triple of $C_c$ is covered by some $U_i$.
\item[(iii)] For every independent pair $F$ of $C_c$ there is an index $i$
such that $U_i$ meets $F$, and every independent triple of $C_c$ that is
covered by $U_i$ and by no other $U_j$ also meets~$F$.
\end{enumerate}
\end{lemma}

\begin{proof}
Restrict the coloring to $\KG(C_c,3)$. By Schrijver's theorem this graph has
chromatic number $c-4$, so the restriction is an optimal coloring that uses
all $c-4$ colors. In an optimal coloring, every color $i$ has a vertex that
has neighbors of all the other colors. Otherwise we could recolor each vertex
of color $i$ with a color missing from its neighborhood. These vertices are
pairwise nonadjacent, so the result is proper, and it uses one color fewer.
Let $U_i$ be such a vertex of color~$i$.

Let $F\subset U_i$ be a pair. The triple $F\cup\{\infty\}$ is independent in
$H$. For every color $j\ne i$, some triple of color $j$ is disjoint from
$U_i$, hence from $F\cup\{\infty\}$. So $F\cup\{\infty\}$ has color~$i$. This
proves~(i), since a pair in $U_i\cap U_j$ would force two colors on the same
vertex.

Let $W$ be an independent triple of $C_c$ of color $i$. If $W$ contains no
pair of $U_i$, then $|W\cap U_i|\le1$, so some pair $F\subset U_i$ is
disjoint from $W$. Then $F\cup\{\infty\}$ and $W$ are disjoint and both have
color $i$, which is impossible. So $W$ is covered by $U_i$. This
proves~(ii), and it shows that a triple covered by $U_i$ alone has color~$i$.

For~(iii), let $F$ be an independent pair of $C_c$ and let $i$ be the color
of $F\cup\{\infty\}$. Every triple of color $i$, including $U_i$ and every
triple covered by $U_i$ alone, meets $F\cup\{\infty\}$. Since these triples
avoid $\infty$, they meet~$F$.
\end{proof}

\begin{lemma}\label{lem:isolate}
For $9\le c\le12$, no family $U_1,\dots,U_{c-4}$ as in
Lemma~\ref{lem:anchors} exists. Hence $C_c\sqcup K_1$ is $3$-tight for
$9\le c\le12$.
\end{lemma}

\begin{proof}
This is an exhaustive computer search, carried out by
\path{checks/isolate/verifier.cpp}. For each $c$ it lists every collection of
$c-4$ independent triples of $C_c$ satisfying (i) and (ii), and tests~(iii) on
each one. The search keeps the set of triples not yet covered. At each step it
picks an uncovered triple $W$ and branches over the independent triples that
cover $W$ and share no pair with the triples already chosen, in a fixed
order. In the branch for a triple $U$, the triples
tried earlier for $W$ are excluded. Every collection satisfying (ii) contains
a first triple, in this order, that covers $W$. So every collection is found
exactly once. Once all triples are covered, the search still lists every way
of completing the collection to $c-4$ triples. There is no symmetry
reduction. The results are as follows.
\begin{center}
\begin{tabular}{rrrrr}
\toprule
$c$ & $|\mathcal I_3(C_c)|$ & search nodes & with (i), (ii) & with (iii)\\
\midrule
$9$ & $30$ & $878$ & $123$ & $0$\\
$10$ & $50$ & $10{,}484$ & $120$ & $0$\\
$11$ & $77$ & $159{,}427$ & $0$ & $0$\\
$12$ & $112$ & $2{,}992{,}235$ & $8$ & $0$\\
\bottomrule
\end{tabular}
\end{center}
No collection satisfies~(iii), so by Lemma~\ref{lem:anchors} no coloring
with $c-4$ colors exists.
\end{proof}

\section{Graphs with at least fourteen vertices}\label{sec:large}

In this section $k=3$ and we show that a minimal counterexample has at most
$13$ vertices. The tool is the classification of maximal intersecting triple
systems by Polcyn and Ruci\'nski~\cite{PR}. For a pair $P$ of vertices let
$\up(P)$ be the set of all triples of the ground set containing~$P$.

\begin{lemma}[{\cite[Theorem~4 and Section~2.4]{PR}}]\label{lem:PR}
Let $n\ge7$ and let $\mathcal F$ be a maximal intersecting family of triples
on an $n$-element set that is not a star. Then there are a set $\mathcal P$
of $b$ pairs and a set $\mathcal X$ of $g$ triples with
\[
\mathcal F=\bigcup_{P\in\mathcal P}\up(P)\cup\mathcal X,
\]
where $(b,g)$ is one of $(3,1)$, $(3,0)$, $(2,3)$, $(1,6)$, or $(0,g)$ with
$g\le10$. In every case $g\le1+3(3-b)$.
\end{lemma}

\begin{proof}
For $n\ge7$, Polcyn and Ruci\'nski show that the maximal intersecting families
of triples are, up to isomorphism, the full star, the families
$H_1(n),\dots,H_{11}(n)$, and the families $K_5\cup(n-5)K_1$, $F_7\cup(n-7)K_1$
and $F_{10}\cup(n-7)K_1$. In their notation $H_1(n)$, $H_2(n)$, $H_3(n)$ and
$H_4(n),H_5(n),H_6(n)$ are written exactly in the displayed form, with
$(b,g)$ equal to $(3,1)$, $(3,0)$, $(2,3)$ and $(1,6)$ respectively. The
remaining families have at most $10$ members in total, and we take
$\mathcal P=\emptyset$ for them.
\end{proof}

For a triple $T\in\mathcal F$ that contains no pair of $\mathcal P$, we say
that $T$ is an \emph{exceptional} member of $\mathcal F$. Thus $\mathcal F$
has at most $g$ exceptional members.

We also need a standard lower bound on the number of triangles.

\begin{lemma}\label{lem:triangles}
A graph $J$ with $n$ vertices and $e$ edges has at least
\[
F_n(e)=\frac{e(4e-n^2)}{3n}
\]
triangles.
\end{lemma}

\begin{proof}
An edge $uv$ lies in at least $d_u+d_v-n$ triangles, since $N(u)$ and $N(v)$
are subsets of an $n$-set. Summing over edges counts each triangle three times,
so $3t(J)\ge\sum_{uv\in E(J)}(d_u+d_v)-ne=\sum_v d_v^2-ne\ge4e^2/n-ne$ by
the Cauchy--Schwarz inequality.
\end{proof}

The function $F_n$ is convex, with $F_n'(e)=(8e-n^2)/(3n)$.

\begin{proposition}\label{prop:large}
Every minimal counterexample at rank $3$ has at most $13$ vertices.
\end{proposition}

\begin{proof}
Let $G$ be a minimal counterexample on $n\ge14$ vertices, and let $m=n-5$. By
Lemma~\ref{lem:minimal}(b), $\KG(G,3)$ has a coloring with $m'\le m$
nonempty color classes, all nonstars. Extend each class to a maximal
intersecting family $\mathcal F_i$ of triples of $V(G)$, where $i=1,\dots,m'$.
A family containing a nonstar is a nonstar, so Lemma~\ref{lem:PR} gives sets
$\mathcal P_i$ of $b_i$ pairs and bounds $g_i\le1+3(3-b_i)$ on the number of
exceptional members of~$\mathcal F_i$. Put
\[
D=\sum_{i=1}^{m'}(3-b_i)\ge0.
\]

Let $J$ be the graph on $V(G)$ whose edges are the pairs that are neither
edges of $G$ nor members of any $\mathcal P_i$. Write $e_0=\binom n2-|E(G)|-3m$. Then
\begin{equation}\label{eq:eJ}
e(J)\ge\binom n2-|E(G)|-\sum_i b_i=\binom n2-|E(G)|-3m'+D\ge e_0+D.
\end{equation}
A triangle $T$ of $J$ is an independent triple of $G$, so it has some color
$i$ and lies in $\mathcal F_i$. It contains no pair of $\mathcal P_i$, so it
is an exceptional member of $\mathcal F_i$. Hence
\begin{equation}\label{eq:tJ}
t(J)\le\sum_i g_i\le m'+3D\le m+3D.
\end{equation}

We now show that $F_n(e_0)>m$ and $F_n'(e_0)>3$. Then $F_n$ is increasing
for $e\ge e_0$, and convexity together with~\eqref{eq:eJ} gives
$t(J)\ge F_n(e(J))\ge F_n(e_0)+DF_n'(e_0)>m+3D$, which
contradicts~\eqref{eq:tJ}. Since $F_n$ and $F_n'$ are increasing for
$e\ge n^2/8$, it suffices to check the two inequalities at a lower bound
for~$e_0$ that is at least $n^2/8$.

If $n=14$, then $m=9$. Two cycles of length at least $9$ need $18$ vertices,
so by Lemma~\ref{lem:minimal}(a) the graph $G$ has a path component, and
$|E(G)|\le n-1=13$. Hence $e_0\ge91-13-27=51$, and
\[
F_{14}(51)=\frac{68}{7}>9,\qquad F_{14}'(51)=\frac{106}{21}>3.
\]

If $n\ge15$, then $|E(G)|\le n$, so
\[
e_0\ge E_n:=\binom n2-n-3(n-5)=\frac{n^2-9n+30}{2}.
\]
We have $E_n\ge4n$, with equality at $n=15$, and $4E_n-n^2=n^2-18n+60>0$.
Therefore
\[
F_n(E_n)=\frac{E_n}{3n}\,(4E_n-n^2)\ge\frac43\,(n^2-18n+60),
\qquad
F_n'(E_n)=\frac{3n^2-36n+120}{3n}.
\]
The first bound exceeds $n-5$ because $4n^2-75n+255>0$ for $n\ge15$; at
$n=15$ it equals $30$, and it is increasing for $n\ge10$. The second exceeds
$3$ because $3n^2-45n+120>0$ for $n\ge15$; at $n=15$ it equals $120$, and it is
increasing for $n\ge8$.
\end{proof}

\section{The remaining small graphs}\label{sec:small}

\begin{proof}[Proof of Theorem~\ref{thm:main}]
The forward implication is Theorem~\ref{thm:necessary} with $k=3$. For the
converse, suppose there is a counterexample at rank $3$, and take a minimal
one. By Proposition~\ref{prop:large} and the remark after
Lemma~\ref{lem:minimal}, we may assume it is $G=C_c\sqcup P_r$ with $c\ge9$,
$r\ge1$ and $n=c+r\le13$. If $r=1$ this contradicts Lemma~\ref{lem:isolate}.
The six remaining graphs are listed in Table~\ref{tab:covers}.

Let $m=n-5$. By Lemma~\ref{lem:minimal}(b), there is a coloring of
$\KG(G,3)$ with at most $m$ colors in which every class is a nonstar. Extend
each class to an inclusion-maximal intersecting subfamily of
$\mathcal I_3(G)$. These extensions are still nonstars, and together they
cover $\mathcal I_3(G)$. It therefore suffices to show the following claim:
\emph{for each graph in Table~\ref{tab:covers}, the independent triples cannot
be covered by $m$ maximal intersecting nonstar subfamilies of
$\mathcal I_3(G)$.}

The maximal intersecting subfamilies of $\mathcal I_3(G)$ are the maximal
cliques of the graph on $\mathcal I_3(G)$ in which two triples are adjacent
when they meet. We list all maximal cliques with the Bron--Kerbosch
algorithm with pivoting~\cite{BK} and keep those with no common vertex
(\path{checks/covers/maximal_families.cpp}). Their numbers are in
Table~\ref{tab:covers}. We then rule out a cover by $m$ of them in two
different ways.

\begin{table}[ht]
\centering
\small
\begin{tabular}{lrrrr}
\toprule
$G$ & $|\mathcal I_3(G)|$ & maximal nonstar families & $m$ & certificates\\
\midrule
$C_9\sqcup P_2$ & $84$ & $1{,}842$ & $6$ & $24$\\
$C_9\sqcup P_3$ & $120$ & $11{,}949$ & $7$ & $26$\\
$C_{10}\sqcup P_2$ & $120$ & $8{,}130$ & $7$ & $16$\\
$C_9\sqcup P_4$ & $165$ & $32{,}871$ & $8$ & $28$\\
$C_{10}\sqcup P_3$ & $165$ & $35{,}974$ & $8$ & $26$\\
$C_{11}\sqcup P_2$ & $165$ & $28{,}864$ & $8$ & $18$\\
\bottomrule
\end{tabular}
\caption{The six graphs $C_c\sqcup P_r$ with $r\ge2$ and $c+r\le13$.}
\label{tab:covers}
\end{table}

\emph{First method: exact search} (\path{checks/covers/exact_cover.cpp}).
The search keeps the set of uncovered triples. It picks an uncovered triple
and branches over the families containing it, excluding in each branch the
families tried earlier for that triple, so every cover is reached. A branch is
cut when some uncovered triple lies in no remaining family, or when the $q$
largest remaining families, for $q$ the number of families still allowed,
cannot cover what is left. The last test is applied both to the number of
uncovered triples and to their total weight under a fixed nonnegative integer
weighting. Both are valid bounds for any such weighting. The search finds no
cover in all six cases.

\emph{Second method: integer certificates.} This is checked by
\path{checks/orbits/verify_orbits.py}. Suppose $\mathcal A_1,\dots,\mathcal A_m$ cover $\mathcal I_3(G)$, with no
family redundant, and let $\mathcal A=\mathcal A_1$ be a largest one, of size
$s$. Then $ms\ge|\mathcal I_3(G)|$. The automorphisms of $G$ (rotations and
reflections of the cycle, and reversal of the path) permute the families, so
we may take $\mathcal A$ from a fixed set of orbit representatives. Let
$\mathcal U=\mathcal I_3(G)\setminus\mathcal A$. Each other family
$\mathcal B$ has size at most $s$. It must cover at least
\[
\max\bigl(1,\ |\mathcal U|-(m-2)s\bigr)
\]
triples of $\mathcal U$: at least one, since it is not redundant, and enough
that the other $m-2$ families can finish the cover. Call such families
\emph{eligible}. For every representative $\mathcal A$ the program checks one
of three certificates: there is no eligible family; some triple of
$\mathcal U$ lies in no eligible family; or there are nonnegative integer
weights on $\mathcal U$ with total weight greater than $(m-1)q$, while every
eligible family has weight at most $q$. Each certificate rules out a cover.
The number of representatives checked is in Table~\ref{tab:covers}. All
arithmetic is in integers.

Both methods confirm the claim. So no minimal counterexample exists, and
Theorem~\ref{thm:main} follows.
\end{proof}

\section{A bound for every rank}\label{sec:allk}

\begin{proof}[Proof of Theorem~\ref{thm:allk}]
Let $k\ge3$ and suppose that a minimal counterexample $G$ has $n\ge k^3-k$
vertices. By Lemma~\ref{lem:minimal}(b), $\mathcal I_k(G)$ is covered by at
most $n-2k+1$ intersecting nonstar families of $k$-sets.

By the Hilton--Milner theorem~\cite{HM}, a nonstar intersecting family of
$k$-subsets of an $n$-set, with $n>2k$, has at most
$\binom{n-1}{k-1}-\binom{n-k-1}{k-1}+1$ members. Here
$\binom{n-1}{k-1}-\binom{n-k-1}{k-1}$ is the number of $k$-sets that contain a
fixed vertex $x$ and meet a fixed $k$-set $A$ not containing $x$. This is the
union of the $k$ families $\up(\{x,a\})$ of $k$-sets containing $x$ and $a$,
for $a\in A$, each of size $\binom{n-2}{k-2}$. Since $k\ge3$, two of them
share a set, so the union has at most $k\binom{n-2}{k-2}-1$ members. Hence
every nonstar intersecting family has at most $k\binom{n-2}{k-2}$ members.

On the other hand, each edge of $G$ lies in $\binom{n-2}{k-2}$ of the
$k$-sets, and $|E(G)|\le n$ by Lemma~\ref{lem:minimal}(a). So
\[
|\mathcal I_k(G)|\ge\binom nk-n\binom{n-2}{k-2}.
\]
Using $\binom nk=\frac{n(n-1)}{k(k-1)}\binom{n-2}{k-2}$, the difference
$|\mathcal I_k(G)|-(n-2k+1)\,k\binom{n-2}{k-2}$ is at least
$\binom{n-2}{k-2}/(k(k-1))$ times
\[
n(n-1)-nk(k-1)-(n-2k+1)k^2(k-1)
=n\bigl(n-1-(k^3-k)\bigr)+k^2(k-1)(2k-1).
\]
At $n=k^3-k$ the right side equals $k(k-1)(2k^2-2k-1)>0$, and it increases
with $n$ for $n\ge k^3-k$. So there are more independent $k$-sets than the
color classes can hold, a contradiction.
\end{proof}

\section*{Verification code}

The programs cited in Sections~\ref{sec:isolate} and~\ref{sec:small} are
available at\\ \url{https://github.com/brandonlign/stable-kneser}. A single
script, \path{checks/run_all.sh}, compiles and runs all of them in under a
minute and needs only a C++17 compiler and Python~3.

\section*{Acknowledgments}
Claude (Anthropic) helped with the proofs, the verification code, and the
manuscript.

\begin{thebibliography}{9}

\bibitem{AFTD}
AFTD, \emph{A cycle plus an isolated vertex is $k$-tight once the cycle has
length at least $k((k-1)^2+1)$}, Open problem OP-193, 12 November 2026,
\url{https://rainwangphy.github.io/aftd-platform/blog/2026-11-12-op-193/}.

\bibitem{AHKS}
C.~Arana, P.~Haxell, T.~Kaiser, and M.~Stehl\'ik, \emph{The chromatic number
of $G$-stable Kneser graphs}, preprint, 2026,
\href{https://arxiv.org/abs/2610.17884}{arXiv:2610.17884}.

\bibitem{BK}
C.~Bron and J.~Kerbosch, \emph{Algorithm 457: finding all cliques of an
undirected graph}, Comm. ACM \textbf{16} (1973), 575--577.

\bibitem{DMM}
H.~Daneshpajouh, F.~Meunier, and G.~Mizrahi, \emph{Colorings of complements
of line graphs}, J. Graph Theory \textbf{98} (2021), 216--233.

\bibitem{HM}
A.~J.~W. Hilton and E.~C. Milner, \emph{Some intersection theorems for systems
of finite sets}, Quart. J. Math. Oxford Ser. (2) \textbf{18} (1967), 369--384.

\bibitem{PR}
J.~Polcyn and A.~Ruci\'nski, \emph{A hierarchy of maximal intersecting triple
systems}, Opuscula Math. \textbf{37} (2017), 597--608.

\bibitem{Schrijver}
A.~Schrijver, \emph{Vertex-critical subgraphs of Kneser graphs}, Nieuw Arch.
Wisk. (3) \textbf{26} (1978), 454--461.

\end{thebibliography}

\end{document}
